\ifdefined\gkver\else\def\gkver{student}\fi
\documentclass[\gkver]{gaokaozhenti}
\begin{document}

\chapter{2024年江西}

\begin{enumerate}
\item
%% number: 1
%% paperName: 2024年江西高考第 1 题 4 分
%% typeId: 1
%% type: 单选题
%% chapter: 电场
%% point: 电势差与电场强度的关系
%% method: 无
%% score: 4
%% degree: 950
%% duplicateId: 0
%% body:
极板间一蜡烛火焰带有正离子、电子以及其他的带电粒子，两极板电压保持不变，当电极板距离减小时，电场强度如何变？电子受力方向？\xzanswer{B}
\onepicture[align=right]{figs/01.png}
\fourchoices[answer=B]
{电场强度增大，方向向左}
{电场强度增大，方向向右}
{电场强度减小，方向向左}
{电场强度减小，方向向右}
%% answer: B
%% memo:
\memoanswer{
由题知，两极板电压保持不变，则根据电势差和电场强度的关系有 $E = \frac{U}{d}$，当电极板距离减小时，电场强度 $E$ 增大，再结合题图可知极板间的电场线水平向左，则可知电子受到的电场力方向向右。

故选 B。
}

\item
%% number: 2
%% paperName: 2024年江西高考第 2 题 4 分
%% typeId: 1
%% type: 单选题
%% chapter: 光学
%% point: 光子说
%% method: 无
%% score: 4
%% degree: 950
%% duplicateId: 0
%% body:
近年来，江西省科学家发明硅衬底氮化镓基系列发光二极管，开创了国际上第三条 LED 技术路线。某氮化镓基 LED 材料的简化能级如图所示，若能级差为 $2.20 \UeV$（约 $3.52\times10^{-19} \UJ$），普朗克常量 $h = 6.63\times10^{-34} \UJcdots$，则发光频率约为\xzanswer{C}
\onepicture[align=right]{figs/02.png}
\fourchoices[answer=C]
{$6.38\times10^{14} \UHz$}
{$5.67\times10^{14} \UHz$}
{$5.31\times10^{14} \UHz$}
{$4.67\times10^{14} \UHz$}
%% answer: C
%% memo:
\memoanswer{
根据题意可知，辐射出的光子能量 $E = 3.52 \times 10^{-19} \UJ$，由光子的能量 $E = h\nu$ 得 $\nu = \frac{E}{h} = 5.31 \times 10^{14} \UHz$。

故选 C。
}

\item
%% number: 3
%% paperName: 2024年江西高考第 3 题 4 分
%% typeId: 1
%% type: 单选题
%% chapter: 直线运动
%% point: 匀变速直线运动
%% method: 无
%% score: 4
%% degree: 850
%% duplicateId: 0
%% body:
某物体位置随时间的关系为 $x = 1 + 2t + 3t^{2}$，则关于其速度与 1 $\Us$ 内的位移大小，下列说法正确的是\xzanswer{C}
\fourchoices[answer=C]
{速度是刻画物体位置变化快慢的物理量，1 $\Us$ 内的位移大小为 6 $\Um$}
{速度是刻画物体位移变化快慢的物理量，1 $\Us$ 内的位移大小为 6 $\Um$}
{速度是刻画物体位置变化快慢的物理量，1 $\Us$ 内的位移大小为 5 $\Um$}
{速度是刻画物体位移变化快慢的物理量，1 $\Us$ 内的位移大小为 5 $\Um$}
%% answer: C
%% memo:
\memoanswer{
根据速度的定义式 $v = \frac{\Delta x}{\Delta t}$ 表明，速度等于位移与时间的比值。位移是物体在一段时间内从一个位置到另一个位置的位置变化量，而时间是这段时间的长度。这个定义强调了速度不仅描述了物体运动的快慢，还描述了物体运动的方向。因此，速度是刻画物体位置变化快慢的物理量。

再根据物体位置随时间的关系 $x = 1 + 2t + 3t^{2}$，可知开始时物体的位置 $x_{0} = 1 \Um$，1 $\Us$ 时物体的位置 $x_{1} = 6 \Um$，则 1 $\Us$ 内物体的位移为 $\Delta x = x_{1} - x_{0} = 5 \Um$。

故选 C。
}

\item
%% number: 4
%% paperName: 2024年江西高考第 4 题 4 分
%% typeId: 1
%% type: 单选题
%% chapter: 曲线运动
%% point: 人造地球卫星
%% method: 无
%% score: 4
%% degree: 650
%% duplicateId: 0
%% body:
两个质量相同的卫星绕月球做匀速圆周运动，半径分别为 $r_{1}$、$r_{2}$，则动能和周期的比值为\xzanswer{A}
\fourchoices[answer=A]
{$\frac{E_{k1}}{E_{k2}} = \frac{r_{2}}{r_{1}}$，$\frac{T_{1}}{T_{2}} = \frac{\sqrt{r_{1}^{3}}}{\sqrt{r_{2}^{3}}}$}
{$\frac{E_{k1}}{E_{k2}} = \frac{r_{1}}{r_{2}}$，$\frac{T_{1}}{T_{2}} = \frac{\sqrt{r_{1}^{3}}}{\sqrt{r_{2}^{3}}}$}
{$\frac{E_{k1}}{E_{k2}} = \frac{r_{2}}{r_{1}}$，$\frac{T_{1}}{T_{2}} = \frac{\sqrt{r_{2}^{3}}}{\sqrt{r_{1}^{3}}}$}
{$\frac{E_{k1}}{E_{k2}} = \frac{r_{1}}{r_{2}}$，$\frac{T_{1}}{T_{2}} = \frac{\sqrt{r_{2}^{3}}}{\sqrt{r_{1}^{3}}}$}
%% answer: A
%% memo:
\memoanswer{
两个质量相同的卫星绕月球做匀速圆周运动，则月球对卫星的万有引力提供向心力，设月球的质量为 $M$，卫星的质量为 $m$，则半径为 $r_{1}$ 的卫星有
$G\frac{Mm}{r_{1}^{2}} = m\frac{v_{1}^{2}}{r_{1}} = m\frac{4\pi^{2}}{T_{1}^{2}}r_{1}$

半径为 $r_{2}$ 的卫星有
$G\frac{Mm}{r_{2}^{2}} = m\frac{v_{2}^{2}}{r_{2}} = m\frac{4\pi^{2}}{T_{2}^{2}}r_{2}$

再根据动能 $E_{k} = \frac{1}{2}mv^{2}$，可得两卫星动能和周期的比值分别为
$\frac{E_{k1}}{E_{k2}} = \frac{r_{2}}{r_{1}}$，$\frac{T_{1}}{T_{2}} = \frac{\sqrt{r_{1}^{3}}}{\sqrt{r_{2}^{3}}}$

故选 A。
}

\item
%% number: 5
%% paperName: 2024年江西高考第 5 题 4 分
%% typeId: 1
%% type: 单选题
%% chapter: 功和能
%% point: 功率
%% method: 无
%% score: 4
%% degree: 650
%% duplicateId: 0
%% body:
庐山瀑布“飞流直下三千尺，疑是银河落九天”瀑布高 150 $\Um$，水流量 10 $\Umcs$，假设利用瀑布来发电，能量转化效率为 70\%，则发电功率为\xzanswer{B}
\fourchoices[answer=B]
{$10^{9}$}
{$10^{7}$}
{$10^{5}$}
{$10^{3}$}
%% answer: B
%% memo:
\memoanswer{
由题知，$\Delta t$ 时间内流出的水量为
$m = \rho Q \Delta t = 1.0 \times 10^{4} \Delta t$

发电过程中水的重力势能转化为电能，则有
$P = \frac{mgh}{\Delta t} \times 70\% \approx 1.1 \times 10^{7} \UW$

故选 B。
}

\item
%% number: 6
%% paperName: 2024年江西高考第 6 题 4 分
%% typeId: 1
%% type: 单选题
%% chapter: 机械波
%% point: 干涉衍射
%% method: 无
%% score: 4
%% degree: 550
%% duplicateId: 0
%% body:
如图~\ref{2024江西06a}所示，利用超声波可以检测飞机机翼内部缺陷。在某次检测实验中，入射波为连续的正弦信号，探头先后探测到机翼表面和缺陷表面的反射信号，分别如图~\ref{2024江西06b}、~\ref{2024江西06c}所示。已知超声波在机翼材料中的波速为 6300 $\Ums$。关于这两个反射信号在探头处的叠加效果和缺陷深度 $d$，下列选项正确的是\xzanswer{A}
\threepicture[labelA=2024江西06a, labelB=2024江西06b, labelC=2024江西06c, wA=4.2cm, wB=3.6cm, wC=3.6cm]
{figs/06a.png}
{figs/06b.png}
{figs/06c.png}
\fourchoices[answer=A]
{振动减弱；$d = 4.725 \Umm$}
{振动加强；$d = 4.725 \Umm$}
{振动减弱；$d = 9.45 \Umm$}
{振动加强；$d = 9.45 \Umm$}
%% answer: A
%% memo:
\memoanswer{
根据反射信号图像可知，超声波的传播周期为 $T = 2 \times 10^{-7} \Us$

又波速 $v = 6300 \Ums$，则超声波在机翼材料中的波长 $\lambda = vT = 1.26 \times 10^{-3} \Um$

结合题图可知，两个反射信号传播到探头处的时间差为 $t = 1.5 \times 10^{-6} \Us$

故两个反射信号的路程差
$2d = v\Delta t = 9.45 \times 10^{-6} \Um = \frac{15}{2}\lambda$

解得 $d = 4.725 \Umm$

两个反射信号在探头处振动减弱，A 正确。

故选 A。
}

\item
%% number: 7
%% paperName: 2024年江西高考第 7 题 4 分
%% typeId: 1
%% type: 单选题
%% chapter: 磁场
%% point: 带电粒子在磁场中的运动
%% method: 无
%% score: 4
%% degree: 450
%% duplicateId: 0
%% body:
石墨烯是一种由碳原子组成的单层二维蜂窝状晶格结构新材料，具有丰富的电学性能。现设计一电路测量某二维石墨烯样品的载流子（电子）浓度。如图~\ref{2024江西07a}所示，在长为 $a$，宽为 $b$ 的石墨烯表面加一垂直向里的匀强磁场，磁感应强度为 $B$，电极 1、3 间通以恒定电流 $I$，电极 2、4 间将产生电压 $U$。当 $I = 1.00\times10^{-3} \UA$ 时，测得 $U - B$ 关系图线如图~\ref{2024江西07b}所示，元电荷 $e = 1.60\times10^{-19} \UC$，则此样品每平方米载流子数最接近\xzanswer{D}
\twopicture[labelA=2024江西07a, labelB=2024江西07b, wA=5.0cm, wB=5.0cm]
{figs/07a.png}
{figs/07b.png}
\fourchoices[answer=D]
{$1.7\times10^{19}$}
{$1.7\times10^{15}$}
{$2.3\times10^{20}$}
{$2.3\times10^{16}$}
%% answer: D
%% memo:
\memoanswer{
设样品每平方米载流子（电子）数为 $n$，电子定向移动的速率为 $v$，则时间 $t$ 内通过样品的电荷量 $q = nevtb$

根据电流的定义式得 $I = \frac{q}{t} = nevb$

当电子稳定通过样品时，其所受电场力与洛伦兹力平衡，则有
$evB = e\frac{U}{b}$

联立解得 $U = \frac{I}{ne}B$

结合图像可得 $k = \frac{I}{ne} = \frac{88\times10^{-3}}{320\times10^{-3}} \UV/\UT$

解得 $n = 2.3 \times 10^{16}$

故选 D。
}

\item
%% number: 8
%% paperName: 2024年江西高考第 8 题 6 分
%% typeId: 2
%% type: 多选题
%% chapter: 曲线运动
%% point: 平抛运动
%% method: 无
%% score: 6
%% degree: 650
%% duplicateId: 0
%% body:
一条河流某处存在高度差，小鱼从低处向上跃出水面，冲到高处。如图所示，以小鱼跃出水面处为坐标原点，$x$ 轴沿水平方向，建立坐标系，小鱼的初速度为 $v_{0}$，末速度 $v$ 沿 $x$ 轴正方向。在此过程中，小鱼可视为质点且只受重力作用。关于小鱼的水平位置 $x$、竖直位置 $y$、水平方向分速度 $v_{x}$ 和竖直方向分速度 $v_{y}$ 与时间 $t$ 的关系，下列图像可能正确的是\xzanswer{AD}
\onepicture[align=right]{figs/08.png}
\fourchoices[answer=AD, ispicture=true, h=2.7cm]
{figs/08a.png}
{figs/08b.png}
{figs/08c.png}
{figs/08d.png}
%% answer: AD
%% memo:
\memoanswer{
AC．小鱼在运动过程中只受重力作用，则小鱼在水平方向上做匀速直线运动，即 $v_{x}$ 为定值，则有水平位移 $x = v_{x}t$

故 A 正确，C 错误；

BD．小鱼在竖直方向上做竖直上抛运动，则
$y = v_{y0}t - \frac{1}{2}gt^{2}$，$v_{y} = v_{y0} - gt$

且最高点时竖直方向的速度为 0，故 B 错误，D 正确。

故选 AD。
}

\item
%% number: 9
%% paperName: 2024年江西高考第 9 题 6 分
%% typeId: 2
%% type: 多选题
%% chapter: 光学
%% point: 双缝干涉测量波长
%% method: 无
%% score: 6
%% degree: 750
%% duplicateId: 0
%% body:
某同学用普通光源进行双缝干涉测光的波长实验。下列说法正确的是\xzanswer{BC}
\fourchoices[answer=BC]
{光具座上依次摆放光源、透镜、滤光片、双缝、单缝、遮光筒、测量头等元件}
{透镜的作用是使光更集中}
{单缝的作用是获得线光源}
{双缝间距越小，测量头中观察到的条纹数目内越多}
%% answer: BC
%% memo:
\memoanswer{
A．进行双缝干涉测光的波长实验，光具座上依次摆放光源、透镜、滤光片、单缝、双缝、遮光筒、测量头等元件，故 A 错误；

B．透镜的作用是使光更集中，故 B 正确；

C．单缝的作用是获得线光源，故 C 正确；

D．根据条纹间距公式 $\Delta x = \frac{L}{d}\lambda$ 可知双缝间距越小，相邻亮条纹的间距较大，测量头中观察到的条纹数目内越少，故 D 错误。

故选 BC。
}

\item
%% number: 10
%% paperName: 2024年江西高考第 10 题 6 分
%% typeId: 2
%% type: 多选题
%% chapter: 电场
%% point: 带电粒子的直线运动
%% method: 无
%% score: 6
%% degree: 450
%% duplicateId: 0
%% body:
如图所示，垂直于水平桌面固定一根轻质绝缘细直杆，质量均为 $m$、带同种电荷的绝缘小球甲和乙穿过直杆，两小球均可视为点电荷，带电荷量分别为 $q$ 和 $Q$。在图示的坐标系中，小球乙静止在坐标原点，初始时刻小球甲从 $x = x_{0}$ 处由静止释放，开始向下运动。甲和乙两点电荷的电势能 $E_{p} = k\frac{Qq}{r}$（$r$ 为两点电荷之间的距离，$k$ 为静电力常量）。最大静摩擦力等于滑动摩擦力 $f$，重力加速度为 $g$。关于小球甲，下列说法正确的是\xzanswer{BD}
\onepicture[align=right]{figs/10.png}
\fourchoices[answer=BD]
{最低点的位置 $x = \frac{kQq}{(mg + f)x_{0}}$}
{速率达到最大值时的位置 $x = \sqrt{\frac{kQq}{mg - f}}$}
{最后停留位置 $x$ 的区间是 $\sqrt{\frac{kQq}{mg}} \leq x \leq \sqrt{\frac{kQq}{mg - f}}$}
{若在最低点能返回，则初始电势能 $E_{p0} < (mg - f)\sqrt{\frac{kQq}{mg + f}}$}
%% answer: BD
%% memo:
\memoanswer{
A．全过程，根据动能定理
$(mg - f)(x_{0} - x) - (k\frac{Qq}{x} - k\frac{Qq}{x_{0}}) = 0$

解得 $x = \frac{kQq}{(mg - f)x_{0}}$

故 A 错误；

B．当小球甲的加速度为零时，速率最大，则有
$mg = f + k\frac{Qq}{x^{2}}$

解得 $x = \sqrt{\frac{kQq}{mg - f}}$

故 B 正确；

C．小球甲最后停留时，满足 $mg - f \leq k\frac{Qq}{x^{2}} \leq mg + f$

解得位置 $x$ 的区间 $\sqrt{\frac{kQq}{mg + f}} \leq x \leq \sqrt{\frac{kQq}{mg - f}}$

故 C 错误；

D．若在最低点能返回，即在最低点满足 $k\frac{Qq}{x^{2}} > mg + f$

结合动能定理
$(mg - f)(x_{0} - x) - (k\frac{Qq}{x} - k\frac{Qq}{x_{0}}) = 0$

又 $E_{p} = k\frac{Qq}{r}$

联立可得 $E_{p0} < (mg - f)\sqrt{\frac{kQq}{mg + f}}$

故 D 正确。

故选 BD。
}

\item
%% number: 11
%% paperName: 2024年江西高考第 11 题 8 分
%% typeId: 4
%% type: 实验题
%% chapter: 运动和力
%% point: 验证牛顿第二定律
%% method: 无
%% score: 8
%% degree: 550
%% duplicateId: 0
%% body:
某小组探究物体加速度与其所受合外力的关系。实验装置如图~\ref{2024江西11a}所示，水平轨道上安装两个光电门，小车上固定一遮光片，细线一端与小车连接，另一端跨过定滑轮挂上钩码。
\begin{enumerate}
	\item
	实验前调节轨道右端滑轮高度，使细线与轨道平行，再适当垫高轨道左端以平衡小车所受摩擦力。
	\item
	小车的质量为 $M_{1} = 320 \Ug$。利用光电门系统测出不同钩码质量 $m$ 时小车加速度 $a$。钩码所受重力记为 $F$，作出 $a - F$ 图像，如图~\ref{2024江西11b}中图线甲所示。
	\item
	由图线甲可知，$F$ 较小时，$a$ 与 $F$ 成正比；$F$ 较大时，$a$ 与 $F$ 不成正比。为了进一步探究，将小车的质量增加至 $M_{2} = 470 \Ug$，重复步骤（2）的测量过程，作出 $a - F$ 图像，如图~\ref{2024江西11b}中图线乙所示。
	\item
	与图线甲相比，图线乙的线性区间\tkanswer[2cm]{}，非线性区间\tkanswer[2cm]{}。再将小车的质量增加至 $M_{3} = 720 \Ug$，重复步骤（2）的测量过程，记录钩码所受重力 $F$ 与小车加速度 $a$，如表所示（表中第 9$\sim$14 组数据未列出）。
	\begin{center}
	\begin{tabular}{|c|c|c|c|c|c|}
	\hline
	序号 & 1 & 2 & 3 & 4 & 5 \\
	\hline
	钩码所受重力 $F/9.8\ \UN$ & 0.020 & 0.040 & 0.060 & 0.080 & 0.100 \\
	\hline
	小车加速度 $a/\Umsq$ & 0.26 & 0.55 & 0.82 & 1.08 & 1.36 \\
	\hline
	序号 & 6 & 7 & 8 & 9$\sim$14 & 15 \\
	\hline
	钩码所受重力 $F/9.8\ \UN$ & 0.120 & 0.140 & 0.160 & $\cdots$ & 0.300 \\
	\hline
	小车加速度 $a/\Umsq$ & 1.67 & 1.95 & 2.20 & $\cdots$ & 3.92 \\
	\hline
	\end{tabular}
	\end{center}
	\item
	请在图~\ref{2024江西11b}中补充描出第 6 至 8 三个数据点，并补充完成图线丙\tkanswer[2cm]{}。
	\item
	根据以上实验结果猜想和推断：小车的质量\tkanswer[2cm]{}时，$a$ 与 $F$ 成正比。结合所学知识对上述推断进行解释：\tkanswer[4cm]{}。
\end{enumerate}
\twopicture[labelA=2024江西11a, labelB=2024江西11b, wA=7.0cm, wB=6.4cm, align=right]
{figs/11a.png}
{figs/11b.png}
%% answer:
\jdanswer{
\begin{enumerate}
	\item
	（实验步骤，无作答）
	\item
	\quad
	\item
	\quad
	\item
	较大，较小
	\item
	见解析图
	\item
	远大于钩码质量。设绳子拉力为 $T$，对钩码根据牛顿第二定律有 $F - T = ma$，对小车根据牛顿第二定律有 $T = Ma$，联立解得 $F = (M + m)a$，变形得 $a = \frac{1}{m + M}F$；当 $m \ll M$ 时，可认为 $m + M = M$，则 $a = \frac{1}{M}F$，即 $a$ 与 $F$ 成正比。
\end{enumerate}
}
%% memo:
\memoanswer{
（4）由题图（b）分析可知，与图线甲相比，图线乙的线性区间较大，非线性区间较小；

（5）在坐标系中进行描点，结合其他点用平滑的曲线拟合，使尽可能多的点在线上，不在线上的点均匀分布在线的两侧，如下图所示
\onepicture[w=7.0cm]{figs/11an.png}

（6）设绳子拉力为 $T$，对钩码根据牛顿第二定律有 $F - T = ma$

对小车根据牛顿第二定律有 $T = Ma$

联立解得 $F = (M + m)a$

变形得 $a = \frac{1}{m + M}F$

当 $m \ll M$ 时，可认为 $m + M = M$

则 $a = \frac{1}{M}F$

即 $a$ 与 $F$ 成正比。
}

\item
%% number: 12
%% paperName: 2024年江西高考第 12 题 7 分
%% typeId: 4
%% type: 实验题
%% chapter: 电路
%% point: 测量电阻率
%% method: 无
%% score: 7
%% degree: 450
%% duplicateId: 0
%% body:
某小组欲设计一种电热水器防触电装置，其原理是：当电热管漏电时，利用自来水自身的电阻，可使漏电电流降至人体安全电流以下。为此，需先测量水的电阻率，再进行合理设计。
\begin{enumerate}
	\item
	如图~\ref{2024江西12a}所示，在绝缘长方体容器左右两侧安装可移动的薄金属板电极，将自来水倒入其中，测得水的截面宽 $d = 0.07 \Um$ 和高 $h = 0.03 \Um$。
	\item
	现有实验器材：电流表（量程 300 $\UmuA$，内阻 $R_{A} = 2500 \UO$）、电压表（量程 3 $\UV$ 或 15 $\UV$，内阻未知）、直流电源（3 $\UV$）、滑动变阻器、开关和导线。请在图~\ref{2024江西12a}中画线完成电路实物连接\tkanswer[2cm]{}。
	\item
	连接好电路，测量 $26^{\circ}$C 的水在不同长度 $l$ 时的电阻值 $R_{x}$。将水温升到 $65^{\circ}$C，重复测量。绘出 $26^{\circ}$C 和 $65^{\circ}$C 水的 $R_{x} - l$ 图，分别如图~\ref{2024江西12b}中甲、乙所示。
	\item
	若 $R_{x} - l$ 图线的斜率为 $k$，则水的电阻率表达 $\rho =$\tkanswer[2cm]{}（用 $k$、$d$、$h$ 表示）。实验结果表明，温度\tkanswer[2cm]{}（填“高”或“低”）的水更容易导电。
	\item
	测出电阻率后，拟将一段塑料水管安装于热水器出水口作为防触电装置。为保证出水量不变，选用内直径为 $8.0\times10^{-3}$ $\Um$ 的水管。若人体的安全电流为 $1.0\times10^{-3} \UA$，热水器出水温度最高为 $65^{\circ}$C，忽略其他电阻的影响（相当于热水器 220 $\UV$ 的工作电压直接加在水管两端），则该水管的长度至少应设计为\tkanswer[2cm]{}$\Um$。（保留两位有效数字）
\end{enumerate}
\twopicture[labelA=2024江西12a, labelB=2024江西12b, wA=7.0cm, wB=6.4cm, align=right]
{figs/12a.png}
{figs/12b.png}
%% answer:
\jdanswer{
\begin{enumerate}
	\item
	（实验步骤，无作答）
	\item
	见解析图
	\item
	\quad
	\item
	$kdh$，高
	\item
	0.46
\end{enumerate}
}
%% memo:
\memoanswer{
（2）电源电动势为 3 $\UV$，故电压表量程选择 3 $\UV$；由于电流表的内阻已知，故采用电流表内接时，可以消除系统误差，故电流表采用内接法，实物图如下：
\onepicture[w=7.0cm]{figs/12an.png}

（4）根据电阻定律 $R_{x} = \rho\frac{l}{dh} = \frac{\rho}{dh}\cdot l$，故可得 $k = \frac{\rho}{dh}$，得 $\rho = kdh$

电阻率越小更容易导电，根据图像可知 $65^{\circ}$C 的水的电阻率更小，故可知温度高的水更容易导电。

（5）根据前面分析可知 $65^{\circ}$C 的水的电阻率为
$\rho = kdh = \frac{8\times10^{3}}{0.7}\times0.07\times0.03 \UOm = 24 \UOm$

故当选用内直径为 $d_{0} = 8.0 \times 10^{-3} \Um$ 的水管。若人体的安全电流为 $I_{0} = 1.0 \times 10^{-3} \UA$，接入电压 $U = 220 \UV$ 时，得
$\frac{U}{R} = I_{0}$，$R = \rho\frac{l_{m}}{\pi\left(\frac{d_{0}}{2}\right)^{2}}$

解得水管的长度至少应设计为 $l_{m} = 0.46 \Um$
}

\item
%% number: 13
%% paperName: 2024年江西高考第 13 题 10 分
%% typeId: 6
%% type: 计算题
%% chapter: 气体性质
%% point: 气体多过程
%% method: 无
%% score: 10
%% degree: 750
%% duplicateId: 0
%% body:
可逆斯特林热机的工作循环如图所示。一定质量的理想气体经 ABCDA 完成循环过程，AB 和 CD 均为等温过程，BC 和 DA 均为等容过程。已知 $T_{1} = 1200 \UK$，$T_{2} = 300 \UK$，气体在状态 A 的压强 $p_{A} = 8.0\times10^{5} \UkPa$，体积 $V_{1} = 1.0 \Umc$，气体在状态 C 的压强 $p_{C} = 1.0\times10^{5} \UkPa$。求：
\begin{enumerate}
	\item
	气体在状态 D 的压强 $p_{D}$；
	\item
	气体在状态 B 的体积 $V_{2}$。
\end{enumerate}
\onepicture[align=right]{figs/13.png}
%% answer:
\jdanswer{
\begin{enumerate}
	\item
	$2\times10^{5} \UkPa$
	\item
	$2.0 \Umc$
\end{enumerate}
}
%% memo:
\memoanswer{
（1）从 D 到 A 状态，根据查理定律
$\frac{p_{A}}{T_{1}} = \frac{p_{D}}{T_{2}}$

解得 $p_{D} = 2\times10^{5} \UkPa$

（2）从 C 到 D 状态，根据玻意耳定律
$p_{C}V_{2} = p_{D}V_{1}$

解得 $V_{2} = 2.0 \Umc$
}

\item
%% number: 14
%% paperName: 2024年江西高考第 14 题 11 分
%% typeId: 6
%% type: 计算题
%% chapter: 曲线运动
%% point: 圆周运动的应用
%% method: 无
%% score: 11
%% degree: 450
%% duplicateId: 0
%% body:
雪地转椅是一种游乐项目，其中心传动装置带动转椅在雪地上滑动。如图~\ref{2024江西14a}、~\ref{2024江西14b}所示，传动装置有一高度可调的水平圆盘，可绕通过中心 O 点的竖直轴匀速转动。圆盘边缘 A 处固定连接一轻绳，轻绳另一端 B 连接转椅（视为质点）。转椅运动稳定后，其角速度与圆盘角速度相等。转椅与雪地之间的动摩擦因数为 $\mu$，重力加速度为 $g$，不计空气阻力。
\begin{enumerate}
	\item
	在图~\ref{2024江西14a}中，若圆盘在水平雪地上以角速度 $\omega_{1}$ 匀速转动，转椅运动稳定后在水平雪地上绕 O 点做半径为 $r_{1}$ 的匀速圆周运动。求 AB 与 OB 之间夹角 $\alpha$ 的正切值。
	\item
	将圆盘升高，如图~\ref{2024江西14b}所示。圆盘匀速转动，转椅运动稳定后在水平雪地上绕 $O_{1}$ 点做半径为 $r_{2}$ 的匀速圆周运动，绳子与竖直方向的夹角为 $\theta$，绳子在水平雪地上的投影 $A_{1}B$ 与 $O_{1}B$ 的夹角为 $\beta$。求此时圆盘的角速度 $\omega_{2}$。
\end{enumerate}
\twopicture[labelA=2024江西14a, labelB=2024江西14b, wA=4.6cm, wB=5.0cm, align=right]
{figs/14a.png}
{figs/14b.png}
%% answer:
\jdanswer{
\begin{enumerate}
	\item
	$\frac{\mu g}{\omega_{1}^{2}r_{1}}$
	\item
	$\omega_{2} = \sqrt{\frac{\mu g\sin\theta\cos\beta}{(\sin\theta\sin\beta + \mu\cos\theta)r_{2}}}$
\end{enumerate}
}
%% memo:
\memoanswer{
（1）转椅做匀速圆周运动，设此时轻绳拉力为 $T$，转椅质量为 $m$，受力分析可知轻绳拉力沿切线方向的分量与转椅受到地面的滑动摩擦力平衡，沿径向方向的分量提供圆周运动的向心力，故可得
$T\cos\alpha = m\omega_{1}^{2}r_{1}$
$\mu mg = T\sin\alpha$

联立解得 $\tan\alpha = \frac{\mu g}{\omega_{1}^{2}r_{1}}$

（2）设此时轻绳拉力为 $T'$，沿 $A_{1}B$ 和垂直 $A_{1}B$ 竖直向上的分力分别为
$T_{1} = T'\sin\theta$，$T_{2} = T'\cos\theta$

对转椅根据牛顿第二定律得
$T_{1}\cos\beta = m\omega_{2}^{2}r_{2}$

沿切线方向 $T_{1}\sin\beta = f = \mu F_{N}$

竖直方向 $F_{N} + T_{2} = mg$

联立解得 $\omega_{2} = \sqrt{\frac{\mu g\sin\theta\cos\beta}{(\sin\theta\sin\beta + \mu\cos\theta)r_{2}}}$
}

\item
%% number: 15
%% paperName: 2024年江西高考第 15 题 18 分
%% typeId: 6
%% type: 计算题
%% chapter: 电磁感应
%% point: 双杆问题
%% method: 无
%% score: 18
%% degree: 350
%% duplicateId: 0
%% body:
如图~\ref{2024江西15a}所示，轨道左侧斜面倾斜角满足 $\sin\theta_{1} = 0.6$，摩擦因数 $\mu_{1} = \frac{3}{20}$，足够长的光滑水平导轨处于磁感应强度为 $B = 0.5 \UT$ 的匀强磁场中，磁场方向竖直向上，右侧斜面导轨倾角满足 $\sin\theta_{2} = 0.8$，摩擦因数 $\mu_{2} = \frac{83}{144}$。现将质量为 $m_{\text{甲}} = 6 \Ukg$ 的导体杆甲从斜面上高 $h = 4 \Um$ 处由静止释放，质量为 $m_{\text{乙}} = 2 \Ukg$ 的导体杆乙静止在水平导轨上，与水平轨道左端的距离为 $d$。已知导轨间距为 $l = 2 \Um$，两杆电阻均为 $R = 1 \UO$，其余电阻不计，不计导体杆通过水平导轨与斜面导轨连接处的能量损失，且若两杆发生碰撞，则为完全非弹性碰撞，取 $g = 10 \Umsq$，求：
\begin{enumerate}
	\item
	甲杆刚进入磁场，乙杆的加速度？
	\item
	乙杆第一次滑上斜面前两杆未相碰，距离 $d$ 满足的条件？
	\item
	若乙前两次在右侧倾斜导轨上相对于水平导轨的竖直高度 $y$ 随时间 $t$ 的变化如图~\ref{2024江西15b}所示（$t_{1}$、$t_{2}$、$t_{3}$、$t_{4}$、$b$ 均为未知量），乙第二次进入右侧倾斜导轨之前与甲发生碰撞，甲在 $0 \sim t_{3}$ 时间内未进入右侧倾斜导轨，求 $d$ 的取值范围。
\end{enumerate}
\twopicture[labelA=2024江西15a, labelB=2024江西15b, wA=6.0cm, wB=4.6cm, align=right]
{figs/15a.png}
{figs/15b.png}
%% answer:
\jdanswer{
\begin{enumerate}
	\item
	$a_{\text{乙}0} = 2 \Umsq$，方向水平向右
	\item
	$d \geq 24 \Um$
	\item
	$\frac{356}{11} \Um < d < \frac{696}{11} \Um$
\end{enumerate}
}
%% memo:
\memoanswer{
（1）甲从静止运动至水平导轨时，根据动能定理有
$m_{1}gh - \mu_{1}m_{1}g\cos\theta_{1}\cdot\frac{h}{\sin\theta_{1}} = \frac{1}{2}m_{1}v_{0}^{2}$

甲刚进入磁场时，平动切割磁感线有 $E_{0} = Blv_{0}$

则根据欧姆定律可知此时回路的感应电流为
$I_{0} = \frac{E_{0}}{2R}$

根据楞次定律可知，回路中的感应电流沿逆时针方向（俯视），结合左手定则可知，乙所受安培力方向水平向右，由牛顿第二定律有
$BI_{0}l = m_{2}a_{\text{乙}0}$

带入数据有
$a_{\text{乙}0} = 2 \Umsq$，方向水平向右

（2）甲和乙在磁场中运动的过程中，系统不受外力作用，则系统动量守恒，若两者共速时恰不相碰，则有
$m_{1}v_{0} = (m_{1} + m_{2})v_{\text{共}}$

对乙根据动量定理有 $B\overline{I}lt = m_{2}v_{\text{共}}$

其中 $\overline{I}t = q = \frac{Bl\Delta x}{2R}$

联立解得 $d_{\min} = \Delta x = 24 \Um$

则 $d$ 满足 $d \geq 24 \Um$

（3）根据（2）问可知，从甲刚进入磁场至甲、乙第一次在水平导轨运动稳定，相对位移为 $\Delta x = 24 \Um$，且稳定时的速度 $v_{\text{共}} = 6 \Ums$

乙第一次在右侧斜轨上向上运动的过程中，根据牛顿第二定律有
$m_{2}g\sin\theta_{2} + \mu_{2}m_{2}g\cos\theta_{2} = m_{2}a_{\text{乙上}}$

根据匀变速直线运动位移与速度的关系有 $2a_{\text{乙上}}x_{\text{上}} = v_{\text{共}}^{2}$

乙第一次在右侧斜轨上向下运动的过程中，根据牛顿第二定律有
$m_{2}g\sin\theta_{2} - \mu_{2}m_{2}g\cos\theta_{2} = m_{2}a_{\text{乙下}}$

再根据匀变速直线运动位移与速度的关系有
$2a_{\text{乙下}}x_{\text{下}} = v_{1}^{2}$

且 $x_{\text{上}} = x_{\text{下}}$

联立解得乙第一次滑下右侧轨道最低点的速度 $v_{1} = 5 \Ums$

由于两棒发生碰撞，则为完全非弹性碰撞，则甲乙整体第一次在右侧倾斜轨道上向上运动有
$(m_{1} + m_{2})g\sin\theta_{2} + \mu_{2}(m_{1} + m_{2})g\cos\theta_{2} = (m_{1} + m_{2})a_{\text{共上}}$

同理有 $2a_{\text{共上}}x_{\text{共上}} = v^{2}$

且由图~\ref{2024江西15b}可知 $x_{\text{上}} = 4.84x_{\text{共上}}$

解得甲、乙碰撞后的速度 $v = \frac{30}{11} \Ums$

乙第一次滑下右侧轨道最低点后与甲相互作用的过程中，甲、乙组成的系统合外力为零，根据动量守恒有
$m_{1}v_{2} - m_{2}v_{1} = (m_{1} + m_{2})v$

解得乙第一次滑下右侧轨道最低点时甲的速度为
$v_{2} = \frac{175}{33} \Ums$

若乙第一次滑下右侧轨道最低点时与甲发生碰撞，则对应 $d$ 的最小值，乙第一次在右侧斜轨上运动的过程，对甲根据动量定理有
$-B\overline{I}_{1}l\Delta t_{1} = m_{1}v_{2} - m_{1}v_{\text{共}}$

其中 $\overline{I}_{1}\Delta t_{1} = q_{1} = \frac{Bl\Delta x_{1}}{2R}$

解得 $\Delta x_{1} = \frac{92}{11} \Um$

根据位移关系有
$d_{\min}' - \Delta x = \Delta x_{1}$

解得 $d_{\min}' = \frac{356}{11} \Um$

若乙返回水平导轨后，当两者共速时恰好碰撞，则对应 $d$ 的最大值，对乙从返回水平导轨到与甲碰撞前瞬间的过程，根据动量定理有
$B\overline{I}_{2}l\Delta t_{2} = m_{2}v + m_{2}v_{1}$

其中 $\overline{I}_{2}\Delta t_{2} = q_{2} = \frac{Bl\Delta x_{2}}{2R}$

解得 $\Delta x_{2} = \frac{340}{11} \Um$

根据位移关系有 $d_{\max} - \Delta x - \Delta x_{1} = \Delta x_{2}$

解得 $d_{\max} = \frac{696}{11} \Um$

则 $d$ 的取值范围为 $\frac{356}{11} \Um < d < \frac{696}{11} \Um$
}

\end{enumerate}

\end{document}
